% !TeX root = ../papers/nc-p3.tex
\section{負付値の二つの曲線近傍と境界指数の比較}
\label{sec:rank-fifteen-negative-x-arcs}
\begin{proposition}\label{prop:rank-fifteen-negative-x-arcs}
負付値座標 $w,g$ がともに正付値であるとする。通常な階数15窓が
\[
   \nu(g-w^2)\ge37\nu(w)\quad\text{または}\quad
   \nu(w-g^2)\ge21\nu(g)
\]
を満たすなら、標準Koszul左右AS正則であり、各始点は四つの中心四次式上で自由階数256である。
初めの近傍は $\nu(y)=0$ に、後の近傍は $2\nu(y)=3\nu(x)<0$ に属する。
\end{proposition}
\begin{proof}
前節の整な倍率表示を $(w,g)=(t,t^2)$ と $(t^2,t)$ に制限する。
平方根の定数項は前節と同じである。
\path{computations/verify_rank_fifteen_negative_x_arcs.py} は
$\mathbb Q(i)$ で展開し、3で還元したノルムが非零の主元を選ぶ。
前者の未切断経路の損失は $(0,2,2,4,4)$、中心は $(0,2,4,4)$、五次の積は四桁を使う。
六次から十三次は $(6,2,2,2,2,0,2,0)$ であり、三十七桁の最後の一桁が残る。
後者の各損失は 順に $(0,1,1,3,2)$、$(0,1,2,2)$、二桁、$(4,1,1,1,1,0,1,0)$ である。
両曲線で全四始点の中心商は十三次で消え、単純根と消去を上記の桁数で比較すると
前節と同様に近傍へ結論が伸びる。$\nu(x)=-4\nu(w)$,
$\nu(y)=4\nu(g)-8\nu(w)$ から水準の付値式も従う。
\end{proof}

単元 $m$ を残す二つの半直線 $g=mw^2$, $w=mg^2$ 上では、それぞれ
$\bar y=-\bar m^4$, $\overline{y^2/x^3}=-1/(64\bar m^4)$ となる。
実際 $xw^4\equiv-4$、$H/g^2\equiv-i$ を $y=x^2H^2/(4+r^2)^2$ に入れればよい。
$(t,t^2)$ の全308組の指数列を非分岐の辺と前節の対角線の項別和と比べると、
一致は254組であり、未切断四次でも16行列のうち4行列が異なる。
$(t^2,t)$ と前節の対角線では 307組が一致する。唯一異なるのは文字1の四次中心で、三つの指数が零から1に増える。最大指数だけを比べるとこの差は見えない。
ここで比較した対角線は倍率 $m=1$ の計算である。一般の半直線で
$\mathbb Q(i)(m)$ 上の指数と例外倍率を求める必要が残る。
